Table~\ref{tab:datastats} presents the statistics of the in‑distribution (ID) and out‑of‑distribution (OOD) datasets; the additional details are as follows:
\subsection{Information for In-Distribution Datasets}
We curate a diverse suite of training tasks integrated into \env, spanning fifteen datasets across eight distinct reasoning domains:
\begin{itemize}
    \item \textbf{Chart Understanding}: {FigureQA}~\cite{ebrahimi2018figureqa} and {ChartQA}~\cite{masry2022chartqa} demand precise quantitative extraction and logical inference over diverse data visualizations, requiring the agent to align visual features (bars, lines) with numerical values.
    \item \textbf{Geometric Reasoning}: {Geometry3K}~\cite{lu2021inter}, {GeoQA}~\cite{chen2021geoqa}, and {UniGeo}~\cite{chen2022unigeo} necessitate grounding symbolic theorem application in diagrammatic parsing to solve multi-step spatial and quantitative problems.
    \item \textbf{Mathematical Reasoning}: {MathVision}~\cite{wang2024measuringmultimodalmathematicalreasoning} and {MathV360}~\cite{shi2024mathllavabootstrappingmathematicalreasoning} evaluate the interpretation of complex mathematical diagrams and symbolic expressions, requiring the synthesis of visual perception with algebraic deduction.
    \item \textbf{Geospatial Reasoning}: {MapQA}~\cite{chang2022mapqa} and {InfographicVQA}~\cite{mathew2022infographicvqa} challenge agents to perform rigorous spatial lookups, legend-symbol alignment, and information retrieval across dense graphical layouts.
    \item \textbf{Visual Scientific Reasoning}: {ThinkVL}~\cite{deng2025openvlthinkercomplexvisionlanguagereasoning}, {VizWiz}~\cite{bigham2010vizwiz}, and {ScienceQA}~\cite{lu2022learn} encompass a broad spectrum of tasks ranging from accessibility-focused visual description to domain-specific multimodal scientific inquiry.
    \item \textbf{Document Understanding}: {DocVQA}~\cite{mathew2021docvqa} tests layout-aware information extraction from unstructured documents, including financial reports, forms, and scanned articles.
    \item \textbf{Spatial Reasoning}: {CLEVR}~\cite{johnson2017clevr} provides a diagnostic environment for compositional logic, assessing the agent's ability to execute complex reasoning chains regarding attribute recognition, spatial relations, and counting in 3D scenes.
    \item \textbf{General Visual Reasoning (Others)}: {A-OKVQA}~\cite{schwenk2022okvqa} requires the retrieval of external world knowledge and commonsense reasoning to address open-ended visual questions.
\end{itemize}

\subsection{Information for Out-of-Distribution Datasets}
To assess generalization, we evaluate on six additional benchmarks featuring distinct reasoning:
\begin{itemize}
    \item \textbf{TABMWP}: TABMWP~\cite{lu2022dynamic} focuses on table-based math word problems requiring cell selection, row/column aggregation, and arithmetic reasoning grounded in semi-structured text, testing symbolic computation beyond visual chart parsing.
    \item \textbf{AI2D}: AI2D~\cite{kembhavi2016diagram} is a diagram-centric science QA benchmark involving annotated figures (labels, arrows, parts); it stresses the joint interpretation of schematic layouts and textual cues via multi-hop reasoning.
    \item \textbf{PlotQA}: PlotQA~\cite{methani2020plotqa} centers on scientific plots (axes, legends, bars/lines) demanding high-precision value reading and aggregation, serving as a complementary but distributionally distinct test bed to ChartQA.
    \item \textbf{CLEVR-Math}: CLEVR-Math~\cite{lindstrom2022clevr} extends compositional reasoning in rendered 3D scenes with arithmetic operations, assessing program-like visual logic coupled with numeric computation.
    \item \textbf{IconQA}: IconQA~\cite{lu2021iconqa} comprises textbook-style diagram and icon questions (often multiple-choice) that emphasize abstract visual relations and symbolic understanding, diverging from natural image statistics.
    \item \textbf{MathVista}: MathVista~\cite{lu2023mathvista} is a holistic mathematical reasoning suite spanning real images, charts, and diagrams; it integrates OCR, quantitative reading, geometry, and multi-step deduction for comprehensive evaluation.
\end{itemize}

\begin{table}[t]
\centering
\renewcommand\arraystretch{0.93}
\caption{Dataset statistics in \env{}.}
\fontsize{8}{10}\selectfont\setlength{\tabcolsep}{0.3em}
\resizebox{1.0\linewidth}{!}{ %
\begin{tabular}{l @{\hskip2.5pt} | c | c | c }
\toprule
\textbf{Dataset} & \textbf{Type} & \textbf{\# of training instances} & \textbf{\# of testing instances}\\
\midrule
FigureQA & ID & 2000 & -- \\
ChartQA & ID & 1500 & 1250 \\
Geometry3k & ID& 2000& 600 \\
GeoQA & ID & 2000 & 500 \\
UniGeo & ID& 2000 & 750 \\
MathVision & ID & 2000 & -- \\
MathV360 & ID & 1500 & -- \\
MapQA & ID & 1500 & 2000 \\
InfographicVQA & ID & 1000 & -- \\
ThinkVL & ID & 2000 & -- \\
VizWiz & ID & 200 & -- \\
ScienceQA & ID & 1000 & -- \\
DocVQA & ID & 970 & -- \\
CLEVR & ID & 1000 & -- \\
A-OKVQA & ID & 1500 & -- \\
\midrule
TABMWP & OOD & -- & 970 \\
AI2D & OOD & -- & 810 \\
PlotQA & OOD & -- &2000 \\
Clever-math & OOD & -- & 930\\
IconQA & OOD & -- & 1500\\
Mathvista & OOD & -- & 1000\\
\midrule
\textbf{Total} & -- &\textbf{ 22,170}&\textbf{12,310}\\
\bottomrule
\end{tabular}
}
\label{tab:datastats}
\end{table}



\section{Toolset Information}
\label{app:tool}

  % "Perception": {
  %   "Detection": ["GroundingDINO", "EasyOCRDetector"],
  %   "Segmentation": ["SAM"],
  %   "OCR": ["EasyOCRRecognizer"]
  % },
  % "ChartUnderstanding": {
  %   "ChartToTable": ["UniChart", "DePlot", "ChartMoe.To\_table", "ChartMoe.extrac\t_data"],
  %   "ChartToSVG": ["OpenCV", "ChartDet", "ChartOCR"],
  %   "ChartToSCG": ["ChartDet", "ChartOCR"],
  %   "CaptioningColor": ["ChartAssistant", "ChartVLM"],
  %   "QAAnalysis": ["ChartMoe.answer", "ChartMoe.analyze", "ChartMoe.describe"]
  % },
  % "DiagramFormalization": {
  %   "CDL": ["DiagramFormalizer.image\_cdl", "DiagramFormalizer.text_cdl", "DiagramFormalizer.construction\_cdl", "DiagramFormalizer.goal\_cdl"],
  %   "SymbolicParsing": ["Inter-GPS", "DiagramFormalizer"]
  % },
  % "Solvers": {
  %   "Multimodal": ["G-LLaVA"],
  %   "Math": ["MultiMath"]
  % }
  
In \env, we enable access to 26 tools from four categories detailed as follows: 

\noindent \textbf{Perception.}
The perception layer provides low‑level visual signals that supply structural cues for higher‑level reasoning, including:
\begin{itemize}
    \item \textbf{Detection}: \texttt{GroundingDINO} is an open-set object detection model that unifies visual and textual modalities, allowing for the detection of arbitrary objects described by natural language. It performs open‑set, text‑conditioned object detection, localizing arbitrary entities referenced by natural‑language queries.
    \item \textbf{Segmentation}: \texttt{SAM} (Segment Anything Model) is a foundational model for promptable image segmentation that can zero-shot segment any object based on flexible inputs like points, boxes, or text. It offers promptable, zero‑shot segmentation, isolating regions given points, boxes, or textual prompts.
    \item \textbf{OCR}: \texttt{EasyOCR} is a versatile and ready-to-use optical character recognition tool supporting over 80 languages and various writing scripts for robust text extraction from images. It conducts multilingual optical character recognition, extracting text from images across diverse scripts and layouts.
\end{itemize}

\noindent \textbf{Chart Understanding.}
The chart layer specializes in plot‑ and table‑centric inference. Chart understanding tool collections like \emph{ChartMoE} use a Mixture‑of‑Experts connector to enhance visual–text alignment, supporting both structured data extraction and high‑level analytical reasoning for chart question answering, including:
\begin{itemize}
    \item \textbf{Chart To Table.}
The \emph{ChartToTable} tool subset converts chart images into structured tabular data suitable for downstream numerical reasoning. 
\texttt{UniChart} and \texttt{DePlot} recover underlying data points by detecting axes, legends, and visual marks, while \texttt{ChartMoE.to\_table} and \texttt{ChartMoE.extract\_data} leverage the ChartMoE backbone to extract series-wise values and metadata (\eg, categories, units) with improved robustness across chart styles.

\item \textbf{Chart To SVG.}
The subset of \emph{ChartToSVG} tools reconstructs vectorized representations of charts. 
\texttt{OpenCV} provides low-level image processing to segment graphical primitives, and \texttt{ChartDet} localizes key chart components (axes, legends, bars, lines, markers); \texttt{ChartOCR} recognizes textual elements (titles, axis labels, tick labels, legends) and anchors them to detected regions, enabling SVG-style reconstruction.

\item \textbf{Chart To SCG.}
The subset of \emph{ChartToSCG} tools maps charts into structured scene graphs (SCG). 
\texttt{ChartDet} and \texttt{ChartOCR} are combined to instantiate nodes for visual and textual elements and edges for relations such as series membership, axis association, and legend–color bindings, yielding a graph representation amenable to symbolic reasoning.

    \item \textbf{Captioning Color.}
The subset of \emph{CaptioningColor} generates natural-language descriptions of charts with explicit grounding in visual encodings. 
\texttt{ChartAssistant} produces concise summaries of chart type, trends, and salient extrema, whereas \texttt{ChartVLM} provides more detailed captions that explicitly reference colors, legends, and value ranges.

    \item \textbf{QA Analysis.}
The \emph{QAAnalysis} supports chart question answering and intermediate analysis. 
\texttt{ChartMoE.answer} directly predicts answers to chart-related questions, while \texttt{ChartMoE.analyze} performs intermediate computations (\eg, differences, ratios, aggregations) over extracted data, and \texttt{ChartMoE.describe} generates explanatory textual analyses that connect the visual structure, quantitative reasoning steps, and final answer.
\end{itemize}

\noindent \textbf{Diagram Formalization.}
This layer converts visual diagrams into symbolic structures amenable to automated reasoning. \emph{DiagramFormalizer} parses geometric and schematic diagrams into ConsCDL and ImgCDL‑style formal languages, enabling symbolic representations for downstream deductive geometric reasoning.
\begin{itemize}
    \item \textbf{CDL.}
The \emph{CDL} tools convert geometric diagrams and problem statements into formal constraint languages suitable for symbolic solvers. 
\texttt{image\_cdl} maps raw diagram images to ImgCDL-style primitives (points, lines, circles, incidences), whereas \texttt{text\_cdl} parses textual problem descriptions into ConsCDL-style constraints and goals. 
\texttt{construction\_cdl} recovers the underlying construction sequence of a diagram, and \texttt{goal\_cdl} extracts explicit target predicates (\eg, lengths, angles, congruence relations), enabling theorem-driven reasoning in downstream geometry solvers.  
    \item \textbf{Symbolic Parsing.}
The \emph{SymbolicParsing} tools build fully symbolic problem representations by combining diagram parsing with formal language construction. 
\texttt{Inter-GPS} parses both the problem text and the associated diagram into a unified formal language and then applies theorem-based symbolic reasoning step by step to solve geometry problems. 
\texttt{DiagramFormalizer} provides complementary parsing endpoints that expose ImgCDL/ConsCDL representations directly, supplying structured inputs for interpretable geometry problem solvers such as Inter-GPS.  
\end{itemize}

\noindent \textbf{Math Solvers.}
The solver layer performs domain‑specific high‑level reasoning on top of structured perceptual inputs:
\begin{itemize}
    \item \textbf{Multimodal}: \texttt{G‑LLaVA} is a domain-specific model that integrates a CLIP-based vision encoder with DeepSeekMath-RL to bridge the gap between visual perception and complex mathematical reasoning. It executes geometry‑aware multimodal reasoning by aligning figure interpretation with textual reasoning through specialized instruction tuning.
    \item \textbf{Math}: \texttt{MultiMath} is a specialized multimodal model for geometric problem solving, optimized through geometric cross-modal alignment and instruction tuning to accurately interpret figures and relationships. It bridges visual perception and symbolic mathematics by combining a CLIP‑based vision encoder with DeepSeekMath‑RL, tackling complex multimodal math problems end‑to‑end.
\end{itemize}

